\documentclass[12pt]{article}
\usepackage{fullpage}
\usepackage{amsmath,amssymb,amsthm}

\newtheorem{theorem}{Theorem}

\newcommand{\Aut}{\operatorname{Aut}}
\newcommand{\spanv}{\operatorname{span}}

\begin{document}
\title{Submission A solution to Problem 1}
\date{}
\maketitle

\begin{center}
{\bf A computably AUT-countable example with no finite existential
automorphism basis}
\end{center}

\begin{theorem}
There is a computably represented countable structure $\mathcal A$ such
that every automorphism of every copy of $\mathcal A$ is computable from
that copy, and hence $\mathcal A$ is computably AUT-countable on a cone,
but for every finite parameter set $P\subseteq A$ there is an
automorphism $\pi$ of $\mathcal A$ whose graph is not
$\Sigma^{in}_1$-definable over $P\cup\pi(P)$.
\end{theorem}

\begin{proof}
Let $V=[\omega]^{<\omega}$, with addition $+$ given by symmetric
difference; thus $V$ is a countably infinite vector space over
$\mathbb F_2$.  Fix a computable sequence $(D_i)_{i\in\omega}$ of finite
subsets of $V$ such that every finite subset of $V$ occurs cofinally
often.

We use a one-sorted computable coding of two sorts.  The domain is the
computable disjoint union
\[
        (\{0\}\times\omega)\ \sqcup\ (\{1\}\times V),
\]
identified effectively with $\omega$.  We write its two unary predicates
as $I$ and $S$.  The language also has a binary relation $<$ and a
ternary relation $R$.  On $I$ the relation $<$ is the usual order of
$\omega$, and $<$ is false if one of its arguments is outside $I$.  For
$i\in I$ and $x,y\in S$, put
\[
        R(i,x,y)\quad\Longleftrightarrow\quad x+y\in D_i,
\]
where we identify elements of the $S$-sort with vectors in $V$; on all
other triples $R$ is false.  This is a computable relational structure.

First compute its automorphism group.  Since $(I,<)$ has order type
$\omega$, every automorphism fixes $I$ pointwise.  If
$f\in\Aut(\mathcal A)$, let $c=f(0_V)$.  For every $h\in V$ choose $i$
with $D_i=\{h\}$.  Then $R(i,x,x+h)$ for all $x\in V$, and preservation
of $R$ gives $f(x+h)=f(x)+h$.  Taking $x=0_V$ gives
\[
        f(h)=c+h  \qquad(h\in V).
\]
Conversely, for each $c\in V$, the map $\tau_c$ which fixes $I$
pointwise and sends $x\in S$ to $x+c$ preserves all differences $x+y$,
hence preserves $R$.  Thus
\[
        \Aut(\mathcal A)=\{\tau_c:c\in V\}.
\]

This already implies computable AUT-countability, with no cone oracle.
Let $\mathcal B\cong\mathcal A$, and let $\sigma\in\Aut(\mathcal B)$.
Transporting $\sigma$ to $\mathcal A$, it is some translation $\tau_c$.
Choose an element $i$ of the $I$-sort of $\mathcal B$ corresponding to
some index with $D_i=\{c\}$.  A Turing program with oracle the atomic
diagram of $\mathcal B$ and with this natural number $i$ hardwired
computes $\sigma$: it queries the sort predicate, outputs $n$ on inputs
$n$ in the $I$-sort, and on inputs $x$ in the $S$-sort it searches for
the unique $y$ with $R(i,x,y)$.  Hence every automorphism of every copy
is computable relative to the copy.

It remains to prove the promised failure of $\Sigma^{in}_1$-definability.
Fix a finite parameter set $P\subseteq A$.  Let
\[
        M=\max(P\cap I)
\]
if $P\cap I\neq\varnothing$, and put $M=-1$ otherwise.  Let
\[
        W=\spanv\Bigl(\bigcup_{i\le M}D_i\Bigr)\le V
\]
(where $W=\{0\}$ if $M=-1$).  The subspace $W$ is finite.  Choose a
nonzero $h\in V\setminus W$, and let $\pi=\tau_h$.  Put
$Q=P\cup\pi(P)$.  Since the quotient $V/W$ is infinite and $Q\cap S$ is
finite, choose $x\in V$ such that the two cosets
\[
        x+W,\qquad x+h+W
\]
are disjoint from $Q\cap S$.  Let $y=x+h$.

We shall show that the graph-pair $(x,y)$ is not isolated inside
$\operatorname{graph}(\pi)$ by any existential formula over $Q$.  This
suffices, because if a countable disjunction of existential formulas
defined exactly the graph, then any disjunct true of $(x,y)$ would itself
have to be contained in the graph.

Let
\[
        \exists \bar z\,\theta(\bar q,u,v,\bar z)
\]
be an existential formula with parameters $\bar q$ from $Q$, and suppose
it is true at $(u,v)=(x,y)$.  Choose witnesses.  Passing to one true
conjunction of literals in a disjunctive normal form of the quantifier
free part, it is enough to preserve the truth values of the finitely many
literals in that conjunction.

Call an $I$-sort value \emph{low} if it is at most $M$, and
\emph{high} otherwise.  All low markers have their associated finite
sets $D_i$ contained in $W$.  Let $F$ be the finite set of all $S$-sort
elements occurring among $x,y$, the parameters, and the chosen witnesses,
and put $C=y+W$.  We shall add a vector $t$ precisely to those members of
$F$ which lie in $C$.  The coset $C$ contains no parameter and does not
contain $x$.

Choose $t\in V$ outside the following finite forbidden set.  First avoid
$t=0$ and all choices causing a shifted element of $F\cap C$ to collide
with an unshifted element of $F\setminus C$.  Also, for every
$i\le M$, every $a\in F\cap C$, every $b\in F\setminus C$, and every
$d\in D_i$, avoid $t=a+b+d$.  These exclusions guarantee that all
equalities and all $R$-literals whose first argument is a low marker are
preserved: within $C$ all differences are unchanged; outside $C$ nothing
is changed; and across the cut the original difference is not in $W$,
while the new difference is kept out of each low $D_i$.

There are also finitely many high $I$-sort equality classes among the
chosen witnesses.  For each such class $\mu$, and for each pair
consisting of a positive and a negative $R$-literal with first argument
in $\mu$ and other two arguments in the $S$-sort, we additionally require
that the two resulting differences after the vector shift be unequal.
If the two differences are shifted in the same way this inequality
already follows from the original true diagram; otherwise it excludes
one further value of $t$.  Thus only finitely many more values are
forbidden, and we choose $t$ avoiding all of them.

Now perform the vector shift by this $t$.  For each high equality class
$\mu$, let $A_\mu$ be the finite set of new differences required by the
positive $R$-literals with first argument in $\mu$, and let $B_\mu$ be
the finite set of new differences forbidden by the negative such
literals.  By the choice of $t$, $A_\mu\cap B_\mu=\varnothing$.  Replace
the high $I$-sort witnesses by new indices, all above $M$ and in the
same equality and order pattern as before, with associated sets exactly
$A_\mu$.  This is possible because each finite subset of $V$ occurs
cofinally often among the $D_i$'s.  Sort predicates, equalities, and
$<$-literals are preserved by construction.  Literals involving $R$ with
an argument in the wrong sort remain false, low-marker $R$-literals were
handled above, and high-marker $R$-literals are made true or false
according as their resulting difference lies in $A_\mu$ or not.

Consequently the same existential formula is true of the new pair
\[
        (x,\;y+t)=(x,\;x+h+t).
\]
But $t\neq0$, so $y+t\neq y=\pi(x)$.  Hence this new pair is not in the
graph of $\pi$.  Therefore no existential formula over
$Q=P\cup\pi(P)$ which is true of $(x,\pi(x))$ can be contained in the
graph of $\pi$, and the graph of $\pi$ is not
$\Sigma^{in}_1$-definable over $Q$.

Since $P$ was arbitrary, the required property holds for every finite
parameter set.
\end{proof}

\end{document}
